% Sección A · Mandato acordado. Máximo 1 página (lo exige el enunciado).
% Requiere en el preámbulo: booktabs, array (los mismos de la sección C).
%
% CIFRAS: todas las de esta sección vienen de la exploración de septiembre y
% deben confirmarse contra output/cifras.md (componente B) antes de entregar:
%   11,2 % · 4,574 y 5,085 USD/L · 0,15–0,19 % de baja en listas
%   3,82 % y 4,40 % de descuento · 22,7 % y 32,5 % de mezcla
%   USD 77,26 y 77,20 por caja · 4,47 / 4,83 / 5,03 USD/L
% Van marcadas solo las dos de titular, para no romper la lectura.
\providecommand{\porconfirmar}[1]{\textbf{#1}\,[por confirmar]}

\section{Mandato acordado}

Cada hallazgo va marcado con su origen: \textbf{(V)} verbalizado por la contraparte, \textbf{(M)} medido por nuestro análisis, \textbf{(I)} inferido por el equipo.

\paragraph{La pregunta.} De las cuatro del README, trabajamos la segunda: el precio promedio por litro que subió, ¿es porque subimos precios o porque vendimos mayor proporción de gran reserva?

\paragraph{Por qué esa.} La elección sale de la reunión del 4 de septiembre. Sin que se lo preguntáramos, la contraparte identificó el origen del encargo: «tenemos este problema de la conversión entre cajas y litros, porque no todas las cajas tienen el mismo número de litros; nos dimos cuenta que teníamos un problema y por eso los llamamos» (V). Nombró el vacío que quiere cerrar, «no se disgrega cepa, reserva, gran reserva ni varietal al tomar la decisión» (V), y declaró que su unidad de medida es el litro (V): de ahí que midamos por litro y no por caja.

El dato confirma que la pregunta no es obvia. Entre el primer y el segundo semestre de 2025 el precio medio por litro de exportación subió \porconfirmar{11,2 \%}, de 4,574 a 5,085 USD/L, mientras los precios de lista bajaron entre 0,15 \% y 0,19 \% en los tres tiers y el descuento medio subió de 3,82 \% a 4,40 \% (M). La participación de gran reserva pasó de 22,7 \% a 32,5 \% de los litros (M). Si el precio efectivo sube mientras las listas bajan, el alza es de mezcla, no de precio (I).

\paragraph{Por qué no las otras.} La primera, sobre el tipo de cambio, se usa como contexto: la contraparte no lo mencionó como problema en toda la reunión y declaró no tener política formal para valorizar exportaciones (V), así que no hay cita que la justifique. La tercera, sobre margen por mercado, no es respondible con lo entregado: la contraparte señaló que faltan los aranceles por tratado de libre comercio, los costos de despacho y el tramo de dólar a euro de los clientes europeos (V); sin ellos, CIF y FOB rinden el mismo margen en nuestros datos, lo que es falso (M). La cuarta, sobre qué fecha usar para convertir, queda como definición secundaria: la contraparte cerró que el año «se cierra por factura, no por despacho» (V), y adoptamos la fecha de factura por consistencia (I).

\paragraph{El hilo de correo.} El hilo se abrió el 8 de septiembre con el acta, donde quedó comprometido confirmar por esa vía la equivalencia de caja a litro. La contraparte ofreció enviar los aranceles por país, los costos de despacho y la información financiera (V).
% ---------------------------------------------------------------------------
% ELEGIR UNA DE LAS DOS FRASES SIGUIENTES Y BORRAR LA OTRA.
%
% (1) Si NO se enviaron las consultas dentro del plazo o no se aceptaron:
No se alcanzó a realizar ninguna consulta dentro del plazo, de modo que la equivalencia de caja a litro, el tratamiento del monto neto y la estacionalidad del alza quedan cerrados como supuestos del equipo (bitácora S1, S3 y S8).
%
% (2) Si sí respondieron, reemplazar por algo de esta forma:
% La contraparte respondió el [fecha] que [lo que dijo, textual] (V), con lo que la definición de caja a litro queda [confirmada / corregida a X litros].
% ---------------------------------------------------------------------------

\paragraph{La definición que la contraparte no cerró.} Es la equivalencia entre caja y litro según el formato, que sostiene toda la métrica. Operaciones afirmó que «no todas las cajas tienen el mismo número de litros» (V). El maestro de costos declara 9,0 litros por caja para los dos únicos formatos, 12$\times$750 y 6$\times$1500 (M), que por aritmética dan 9 litros exactos. Y el precio medio por caja del magnum, USD 77,26, coincide con el del gran reserva estándar, USD 77,20 (M), lo que solo se sostiene si ambos rinden los mismos litros (I).

Ante la contradicción, la cerramos en 9,0 litros por caja para ambos formatos. El criterio: cuando lo que dice el cliente y lo que muestra el dato no coinciden, adoptamos el dato si existe una verificación independiente de él. Aquí la hay, y viene de otro archivo: los litros salen del maestro de costos y el precio por caja de la base de ventas. Queda en la bitácora (S1) con su efecto: si el magnum no rindiera 9 litros, el precio medio se movería entre \porconfirmar{4,47 y 5,03 USD/L} frente a los 4,83 adoptados (M); la sensibilidad completa va en la sección D.
